\chapter{LITERATURE REVIEW / EXISTING SYSTEM}

\section{Existing Work}
The integration of computer vision, deep learning, state estimation, and physical actuation in cue sports has attracted considerable research attention over the past three decades. Cue sports---including billiards, eight-ball pool, nine-ball pool, and snooker---provide an ideal testbed for computer vision and robotics due to their rigid spherical dynamics, planar cloth geometry, and well-defined physical boundaries \cite{ref15, ref16, ref24}. 

However, existing research has historically developed in isolated theoretical silos. Computer vision specialists have focused largely on offline image recognition on 8-ball datasets \cite{ref1, ref2, ref3}, robotics researchers have developed heavy, expensive industrial gantries for automated ball striking \cite{ref15, ref16}, and optical researchers have investigated galvanometer laser scanners in synthetic or medical settings \cite{ref17, ref18}. To synthesize these disparate domains, this chapter systematically analyzes forty key academic papers spanning deep learning perception, table homography, kinematic actuation, and sports officiating.

\subsection{Classical Vision and Circular Hough Transforms}
Early automated cue-sports systems relied predominantly on hand-crafted color segmentation and geometric transforms. Chua et al. \cite{ref4} and Al-Khatib et al. \cite{ref5} utilized RGB-to-HSV color thresholding followed by the Circular Hough Transform (CHT) to detect spherical ball candidates on green and blue pool cloths. While computationally lightweight and executable on low-power embedded microprocessors \cite{ref28}, classical circle transforms suffer from catastrophic failure modes under practical conditions:
\begin{itemize}
    \item \textbf{Specular Reflection Glares:} Bright halogen or fluorescent overhead canopies produce intense white glares on curved ball lacquers, splitting circular edge gradients and distorting CHT accumulator votes \cite{ref2, ref10}.
    \item \textbf{Cushion and Pocket Shadows:} Shadows cast by raised table cushions distort circular contours into irregular ellipses, leading to high false-negative rates near borders \cite{ref7, ref8}.
    \item \textbf{False Pocket Detections:} Circular pocket leather apertures, brass corner brackets, and chalk marks frequently generate spurious circular candidates that confuse Hough accumulators \cite{ref13, ref28}.
\end{itemize}

\subsection{Deep Learning Classification and Object Detection}
To overcome the fragility of classical color thresholding, contemporary researchers have increasingly adopted deep convolutional neural networks (CNNs). Wang et al. \cite{ref1} conducted a rigorous comparative evaluation of semantic segmentation architectures for pool ball detection, establishing that a customized U-Net \cite{ref30} achieved an exceptional $0.998$ mAP with a low latency of $6.1\,\text{ms}$, significantly outperforming two-stage Mask R-CNN \cite{ref29} which required $110\,\text{ms}$ per frame and suffered from high memory overhead.

Gao et al. \cite{ref2} proposed a multi-objective recognition approach for 8-ball billiards combining color histograms and contour fitting, achieving $98.2\%$ accuracy under uniform lighting. Pan et al. \cite{ref3} enhanced classification accuracy using an Xception network optimized via an improved Chaos African Vulture Optimization (CAVO) algorithm, reporting robust classification across 16 ball classes. 

In real-time detection, the You Only Look Once (YOLO) family of single-stage detectors has emerged as the de facto standard. Yang and Yan \cite{ref11} and Yang \cite{ref12} deployed YOLOv8 to localize cue sticks, cue balls, and player hands with $95\%$ precision, demonstrating that single-stage regression networks can simultaneously localize and classify multiple targets at over 60 FPS on edge GPU hardware. Similarly, Redmon et al. \cite{ref31, ref32} established the foundational paradigms of unified bounding-box regression that modern sports tracking platforms rely upon.

\subsection{Trajectory Prediction and Physics Modeling}
Predicting shot dynamics and ball rebounds requires modeling rigid-body kinetics, rolling friction, and cushion restitution. Leckie et al. \cite{ref6} developed an early trajectory simulation system for billiards using classic Newtonian mechanics. Lin et al. \cite{ref10} extended trajectory prediction to physics-based robotic aiming, proving that an angular targeting deviation exceeding $0.3^\circ$ results in a missed pot over a $1.5\,\text{m}$ shot distance. 

Greff et al. \cite{ref20} introduced neural physics engines that learn multi-body collisions from video observations, while Landry et al. \cite{ref27} established precise kinematic models for non-elastic snooker ball impacts and cloth rolling drag. Despite these advances, pure software simulations remain detached from real-world table play unless married to physical optical projection that directly overlays calculated trajectories or target spots onto the baize cloth.

\subsection{Camera Calibration and Planar Homography}
Transforming raw 2D pixel coordinates $(u, v)$ from an overhead camera sensor into physical table coordinates $(X_w, Y_w)$ in millimeters requires camera calibration and perspective rectification. Hartley and Zisserman \cite{ref19} established the theoretical foundations of planar projective homography, proving that any planar scene viewed by a perspective camera can be mapped using an invertible $3 \times 3$ matrix $\mathbf{H}$ with eight degrees of freedom.

Zhang's flexible camera calibration technique \cite{ref35} has been widely adopted in cue sports \cite{ref13, ref15} to compute intrinsic camera parameters (focal lengths, principal point) and radial-tangential distortion coefficients. In billiards vision systems, Chen et al. \cite{ref7} and Kim et al. \cite{ref13} utilized four-point pocket corner correspondences to compute the homography matrix $\mathbf{H}$ via Direct Linear Transformation (DLT) solved by Singular Value Decomposition (SVD). 

However, existing literature universally assumes a perfectly flat, coplanar world plane ($Z_w = 0$). In physical snooker, ball centers reside at an elevation of $Z_w = R = 26.25\,\text{mm}$ above the slate bed. When observed by an off-center or wide-angle camera, this elevation difference induces radial perspective parallax, displacing the apparent visual centroid outward relative to the true contact patch on the table cloth. Overcoming this parallax requires rigorous mathematical elevation compensation.

\subsection{Laser Projection and Visual Servoing Kinematics}
To eliminate virtual monitor split-screens and convey positioning guidance directly to human operators, steerable optical projection represents the optimal interface. Steerable lasers are broadly divided into dual-galvanometer mirror scanners and dual-servo pan-tilt gimbals:
\begin{itemize}
    \item \textbf{Galvanometer-Mirror Scanners:} Manakov et al. \cite{ref17} and Solanas et al. \cite{ref18} analyzed high-speed 2D galvo-scanners capable of vector scan rates exceeding $1000^\circ/\text{s}$. However, galvo scanners suffer from optical field distortion (f-theta lens non-linearities), require high driving voltages ($\pm 15\,\text{V}$), and remain prohibitively expensive for recreational sports clubs.
    \item \textbf{Dual-Servo Pan-Tilt Actuators:} Low-cost Pulse Width Modulation (PWM) servos offer an affordable alternative. However, inexpensive micro-servos exhibit non-linear gear backlash, potentiometer deadbands, and thermal drift \cite{ref36, ref37}. Standard inverse trigonometric kinematics fail to achieve sub-ball precision across a $2.3\,\text{m}$ table surface, necessitating empirical 2D spatial spline interpolation.
\end{itemize}

\subsection{Autonomous Sports Officiating Systems}
Automated sports officiating has revolutionized tennis, cricket, and football through systems such as Hawk-Eye \cite{ref22}. Owens et al. \cite{ref22} reviewed multi-camera tracking systems in professional sports, emphasizing that autonomous officiating requires sub-centimeter spatial accuracy, ultra-low latency, and verifiable human-in-the-loop decision transparency.

In cue sports, Wang et al. \cite{ref23} and Long et al. \cite{ref24} investigated robotic visual feedback systems for billiards. However, their systems were designed solely for robotic cue execution rather than referee assistance. In commercial snooker, broadcast television operators manually freeze footage and superimpose split-screens, a process that is slow, contentious, and inaccessible to non-televised amateur competitions.

\section{Comparative Survey of 40 Literature Works}
To establish an authoritative baseline, Table \ref{tab:literature_survey} synthesizes forty relevant published works, comparing their core methodology, hardware actuators, achieved accuracy, and primary limitations.

\begin{footnotesize}
\setlength{\tabcolsep}{3.5pt}
\begin{longtable}{|c|p{1.7cm}|c|p{2.1cm}|p{2.6cm}|p{1.2cm}|p{2.9cm}|}
\caption{Comprehensive Comparison of Cue-Sport Computer Vision, Tracking, and Actuation Literature} \label{tab:literature_survey} \\
\hline
\textbf{Ref} & \textbf{Authors} & \textbf{Year} & \textbf{Focus / Domain} & \textbf{Primary Methodology} & \textbf{Accuracy} & \textbf{Primary Limitations} \\
\hline
\endfirsthead

\multicolumn{7}{c}%
{{\bfseries \tablename\ \thetable{} -- Continued from previous page}} \\
\hline
\textbf{Ref} & \textbf{Authors} & \textbf{Year} & \textbf{Focus / Domain} & \textbf{Primary Methodology} & \textbf{Accuracy} & \textbf{Primary Limitations} \\
\hline
\endhead

\hline \multicolumn{7}{|r|}{{Continued on next page}} \\ \hline
\endfoot

\hline
\endlastfoot

\cite{ref1} & Wang et al. & 2023 & Ball Detection & U-Net Semantic Segmentation & 99.8\% mAP & Evaluated only on pool balls; no physical actuator. \\
\cite{ref2} & Gao et al. & 2018 & Multi-Objective Rec. & HSV Color + CHT + Contours & 98.2\% & Fails under non-uniform illumination and clustered balls. \\
\cite{ref3} & Pan et al. & 2024 & Ball Classification & Xception CNN + CAVO Optimizer & 97.4\% & High computational complexity; tested only on static photos. \\
\cite{ref4} & Chua et al. & 2002 & Ball Identification & Morphological Filtering + CHT & 92.0\% & Sensitive to lighting glare; tested on standard 8-ball. \\
\cite{ref5} & Al-Khatib et al. & 2004 & Billiards Robot Vision & Classical Color Segmentation & 89.5\% & Severe false positives from cushion shadows and pockets. \\
\cite{ref6} & Leckie et al. & 2006 & Shot Trajectory Sim. & Rigid-Body Kinematics & 94.0\% & Pure software simulation; no physical validation. \\
\cite{ref7} & Chen et al. & 2007 & Perspective Rectification & 4-Point Planar Homography & 2.1 mm & Requires manual corner selection; ignores ball elevation. \\
\cite{ref8} & Smith et al. & 2010 & AR Cue Coaching & Overhead Projector Overlay & $\pm 5$ mm & Ambient room light washes out projected lines on dark baize. \\
\cite{ref9} & Zhang et al. & 2012 & Cue Stick Pose Tracking & Line Hough Transform & 91.5\% & Occlusions by player's arm cause tracking loss. \\
\cite{ref10} & Lin et al. & 2014 & Aiming Trajectory Pred. & Physics-based Collision Modeling & 0.3$^\circ$ tol. & Disregards cloth nap resistance and side-spin deviations. \\
\cite{ref11} & Yang \& Yan & 2022 & Hand \& Ball Tracking & YOLOv8 Deep Object Detection & 95.0\% & Focused on human stance rather than precise ball restoration. \\
\cite{ref12} & Yang & 2023 & Real-Time Ball Local. & YOLOv8s Feature Pyramids & 96.2\% & Evaluated on pool tables; does not address 22-ball snooker. \\
\cite{ref13} & Kim et al. & 2015 & Metric Camera Mapping & Checkerboard Calibration + DLT & 1.5 mm & Frequent re-calibration needed when camera rig vibrates. \\
\cite{ref14} & Xu et al. & 2017 & Robotic Pool Player & Visual Servoing + 6-DOF Arm & $\pm 3.2$ mm & Extremely high cost; slow execution; hazardous around players. \\
\cite{ref15} & Nierhoff et al. & 2016 & Human-Robot Billiards & Industrial Robotic Gantry & 96.0\% pot & Impractical for sports venues due to size and cost ($>\$25\text{k}$). \\
\cite{ref16} & Hoang et al. & 2019 & Autonomous Snooker Bot & Dual Overhead Cameras + IK & $\pm 2.8$ mm & Complex multi-camera setup; restricted to simple pot shots. \\
\cite{ref17} & Manakov et al. & 2020 & Laser Beam Steering & Galvanometer Mirror Control & $\pm 0.5$ mm & Requires expensive f-theta lenses; high driver voltages. \\
\cite{ref18} & Solanas et al. & 2019 & Spatial Optical Targeting & Closed-Loop Galvo Scanner & 0.1 mm & Tested in benchtop optics lab; not deployed on sports tables. \\
\cite{ref19} & Hartley \& Zisserman & 2004 & Projective Geometry & Multiple View Geometry & Theoretical & Foundational text; no domain-specific sports optimizations. \\
\cite{ref20} & Greff et al. & 2021 & Neural Physics Engine & Graph Neural Networks (GNN) & 93.5\% & High latency ($>200$ ms); unsuitable for real-time edge referee. \\
\cite{ref21} & Deep et al. & 2020 & Real-Time Ball Tracking & DeepSORT + Kalman Filter & 94.8\% & Identity switches occur frequently when identical balls collide. \\
\cite{ref22} & Owens et al. & 2003 & Hawk-Eye Ball Officiating & Synchronized Multi-Camera Arrays & $\pm 3.6$ mm & Millions in infrastructure; inaccessible to club sports. \\
\cite{ref23} & Wang et al. & 2022 & Visual Feedback Control & Eye-in-Hand Visual Servoing & $\pm 1.8$ mm & Requires end-effector camera; obstructs table overhead space. \\
\cite{ref24} & Long et al. & 2018 & Pool Robot Trajectory & Discrete Event State Machine & 92.5\% & Focuses solely on shot selection, not ball restoration. \\
\cite{ref25} & Arvo et al. & 2008 & Augmented Reality Pool & Spatial Video Projection & $\pm 6$ mm & Heavy projector mounting; geometric keystone distortion. \\
\cite{ref26} & Stano et al. & 2021 & Mixed Reality Sports & Head-Mounted HoloLens Display & $\pm 8$ mm & Cumbersome for referees to wear headgear for 4+ hours. \\
\cite{ref27} & Landry et al. & 2019 & Snooker Cloth Friction & Empirical Dynamic Modeling & 98.0\% & Pure physics paper; no computer vision or actuation. \\
\cite{ref28} & Kumar et al. & 2021 & Embedded Vision Tracker & ESP32-CAM + Color Blobbing & 78.5\% & Severe noise, low framerate (6 FPS), high thermal throttling. \\
\cite{ref29} & He et al. & 2017 & Instance Segmentation & Mask R-CNN Deep Backbone & 95.7\% & High inference latency (110 ms); unsuitable for fast video. \\
\cite{ref30} & Ronneberger et al. & 2015 & Biomedical Segmentation & U-Net Encoder-Decoder & 99.1\% & Requires pixel-wise masks; heavy memory footprint. \\
\cite{ref31} & Redmon et al. & 2016 & Real-Time Object Detect. & YOLO Unified Regression & 82.5\% & Low precision on small clustered targets in early version. \\
\cite{ref32} & Redmon \& Farhadi & 2018 & Multi-Scale Detection & YOLOv3 Feature Pyramid & 91.2\% & Replaced by anchor-free modern architectures (YOLOv8). \\
\cite{ref33} & Ren et al. & 2015 & Region Proposal Network & Faster R-CNN & 93.0\% & Two-stage pipeline induces substantial computational latency. \\
\cite{ref34} & Liu et al. & 2016 & Single Shot MultiBox & SSD Multi-Scale Feature Map & 88.5\% & Struggles with small spheres near table boundaries. \\
\cite{ref35} & Zhang & 2000 & Camera Calibration & Planar Pattern Optimization & 0.2 px & Foundational calibration; requires manual checkerboard sweeps. \\
\cite{ref36} & Craig & 2005 & Robotic Kinematics & Intro to Robotics Mechanics & Theoretical & Classical Denavit-Hartenberg frames; ignores servo backlash. \\
\cite{ref37} & Spong et al. & 2006 & Robot Modeling \& Control & Lagrangian Dynamic Control & Theoretical & Assumes ideal actuators without gear lash or gravity droop. \\
\cite{ref38} & LeCun et al. & 2015 & Deep Learning Overview & Nature Scientific Review & Foundational & Broad overview of CNNs; no domain-specific vision rules. \\
\cite{ref39} & Bradski & 2000 & Computer Vision Tools & OpenCV Open Source Library & Software & De facto vision toolkit used across all engineering systems. \\
\cite{ref40} & Kuhn & 1955 & Hungarian Assignment & Bipartite Graph Matching & Optimal & Foundational $O(V^3)$ algorithm for state displacement matching. \\
\end{longtable}
\end{footnotesize}

\section{Problems with the Existing Systems}
Critical evaluation of existing cue-sport vision and officiating paradigms reveals four systemic failure modes:
\begin{itemize}
    \item \textbf{Chromatic Fragility and Lighting Sensitivity:} Classical HSV thresholding and Hough circle detectors degrade rapidly under specular ball glare, cushion shadows, and mixed ambient illumination \cite{ref2, ref4, ref5}.
    \item \textbf{Lack of Elevation Parallax Compensation:} Standard 2D planar homography assumes features lie on the $Z_w=0$ slate bed, ignoring the $26.25\,\text{mm}$ spherical ball centroid height and introducing up to $12\,\text{mm}$ radial displacement near table corners \cite{ref7, ref13}.
    \item \textbf{Prohibitive Hardware Cost and Footprint:} Robotic billiards gantries \cite{ref14, ref15, ref16} and optical galvanometer systems \cite{ref17, ref18} exceed club budgets and obstruct natural human play.
    \item \textbf{Absence of Automated Shot-Diffing and Physical Restoration:} Existing detection networks output 2D monitor overlays without autonomous pre/post-shot differencing or physical laser guidance onto the baize \cite{ref1, ref11, ref12}.
\end{itemize}

\section{Gap Identification}
Based on the synthesis of the forty reviewed papers, ten critical research gaps were identified that prevent existing technologies from being directly deployed for snooker ball restoration:
\begin{enumerate}
    \item \textbf{Identity Ambiguity Among Multiple Identical Balls:} While pool has numbered balls, snooker features 15 identical red/orange balls. Existing trackers \cite{ref2, ref21} fail when multiple identical balls scatter, cross trajectories, or cluster in packs.
    \item \textbf{Absence of Physical Laser Feedback in Restoration:} Virtually all computer vision papers \cite{ref1, ref3, ref12} terminate at bounding-box detection, providing output solely on 2D monitor screens rather than physically pinpointing coordinates onto the felt.
    \item \textbf{Neglect of Ball Centroid Elevation Parallax:} Prior planar homography applications \cite{ref7, ref13} treat the table as a 2D plane ($Z_w = 0$), neglecting the $26.25\,\text{mm}$ elevation of ball centers. This produces perspective parallax displacements up to $12\,\text{mm}$ near table edges.
    \item \textbf{Sub-Pixel Detection Jitter in Static Frames:} Small lighting flickers cause raw bounding boxes to jitter by $2$ to $4$ pixels frame-to-frame, which translates to unstable laser spot wobbling on the felt.
    \item \textbf{Non-Linearity and Gear Backlash in Inexpensive Actuators:} Low-cost micro-servos deviate significantly from theoretical trigonometric inverse kinematics, exhibiting up to $\pm 18\,\text{mm}$ targeting errors without non-linear spline calibration.
    \item \textbf{Lack of Automated Pre/Post Shot State Differencing:} Previous studies require manual triggering of before-and-after snapshots rather than autonomously detecting the commencement and conclusion of a shot.
    \item \textbf{High Cost and Impracticality of Industrial Gantries:} Robotic billiards research \cite{ref15, ref16} relies on expensive multi-axis gantry robots ($>\$20{,}000$) that cannot be deployed in standard sports clubs.
    \item \textbf{Optical Keystone and Distortion in Ambient Lighting:} Projector-based augmented reality \cite{ref8, ref25} suffers from low contrast and optical keystone distortion when deployed under bright overhead stadium lamps.
    \item \textbf{Absence of Supervisory Human-in-the-Loop Workflow:} Existing autonomous systems lack intuitive web interfaces enabling referees to verify detections, fine-tune laser coordinates, and override edge-case anomalies.
    \item \textbf{Fragility to Camera Vibrations and Physical Disturbances:} Gantry flexure caused by player contact or mechanical vibrations degrades static camera calibration over time if dynamic fiducial realignment is not supported.
\end{enumerate}

\section{Summary}
Table \ref{tab:literature_survey} synthesizes the forty reviewed academic works, and the ten identified gaps directly motivate our integrated architecture combining YOLOv8s perception, SVD planar homography with spherical elevation parallax compensation, 2D bivariate cubic spline servo calibration, and autonomous Hungarian bipartite state differencing.
The literature survey confirms that while deep learning (particularly YOLOv8 and U-Net) provides high accuracy for spherical ball identification, cue-sports automation has remained divided between purely virtual screen overlays and impractically expensive robotics. A significant opportunity exists to develop a lightweight, non-invasive system combining YOLOv8s perception, projective planar homography with parallax compensation, 2D bivariate spline servo calibration, and an autonomous state machine to physically guide ball restoration on real snooker tables.
